% !TEX root = ../main.tex
\chapter{Differentiable manifolds}\label{ch:c03}
\begin{paracol}{2}

\begin{sources}
\begin{tabularx}{\linewidth}{@{}l l L@{}}
\toprule
\textbf{Video} & \textbf{Length} & \textbf{Content}\\
\midrule
\seg{3}{1}{Lecture 3: Differentiable manifolds (Differenzierbare Mannigfaltigkeiten)} & 1:45:19 & Motivation (velocity vectors need
differentiation); two views on a topological manifold (manifold level versus chart level); commutative diagrams;
chart representation of a curve; atlas-continuity and the lemma that makes it equivalent to continuity; chart independence;
chart transition maps and $C^k$-compatibility; $C^k$ atlases, maximal atlases, $C^k$ manifolds; differentiability of
functions through charts and well-definedness (needs $k\ge m$); terminology (smooth manifold); the worked example of $\R$ with
the charts $x(m)=m$ and $y(m)=m^3$.\\
\bottomrule
\end{tabularx}

\smallskip
\textbf{Transcript} \file{materials/transcripts/C03-V01.txt} (German YouTube auto-captions).
\end{sources}

\begin{reading}[Where this lecture sits in the other notes]
\emph{Geometric Anatomy}: lecture 6 (topological manifolds, charts, atlases, transition maps) and lecture 7
(differentiable structures: $C^k$ compatibility, maximal atlases, smooth manifolds, $C^k$ maps between manifolds).
\emph{Gravity and Light} (TP-1 version): lecture 4 (differentiable manifolds).
\end{reading}

% ==================================================================
\section{Why differentiable manifolds?}\label{sec:dm-motivation}

Last time: topological spaces and topological manifolds, introduced because they have just enough structure to define
what a continuous curve or a continuous function is, a motivation deeply rooted in physics \ts{3}{1}{0:00}.
A topological manifold is a very special topological space, one that looks around every point like a piece of $\R^d$;
this lets us say that physical space or spacetime looks everywhere like a piece of $\R^d$ \emph{without} being
globally $\R^d$: it may be a sphere in two dimensions, a three-sphere, or something more complicated; this is the basis of
Hamiltonian mechanics, celestial mechanics, and general relativity \ts{3}{1}{1:06}.

What is new today is again deeply anchored in our idea of physical space \ts{3}{1}{3:16}: when we look at the
track of a particle we want to define at every point of the curve a \emph{velocity vector} (we speak of velocities and
later of accelerations). To extract a velocity vector from a track one must differentiate, and on a topological manifold
there is not enough structure to differentiate; hence we introduce \textbf{differentiable manifolds}: topological manifolds
with additional structure that allows us to differentiate curves, functions, maps in general
\ts{3}{1}{3:36}. Both continuity and differentiability are implicit
assumptions of classical physics \ts{3}{1}{5:23}.
\begin{coursenote}[Velocity, acceleration, and the second derivative]
The first derivative of the track is the tangent vector, the velocity vector. ``That the second derivative of the track
is the acceleration is not quite true; we will come to it''; but we must be able to differentiate at least twice
\ts{3}{1}{5:44}. (On problem sheet~4, after the Easter holidays, one learns about the confusions around the
first and second derivative of a track.) \ts{3}{1}{3:57}
\end{coursenote}

% ==================================================================
\section{Two views on a topological manifold}\label{sec:dm-two-views}

``We have seen this last time, and what I do in this section is a pure \emph{pedagogical trick}: you do not need two
views on a topological manifold, only one. Still, the second view, which you do not need for topological manifolds,
leads to how we proceed with differentiable manifolds'' \ts{3}{1}{6:26}. There are two
completely equivalent ways to decide continuity of objects on a topological manifold \ts{3}{1}{7:09}.

\subsection{View A: the manifold level}\label{ssec:dm-view-a}

Take a curve $\gamma\colon\R\to M$ (parameter range the whole real line) into the manifold \ts{3}{1}{7:52}.
It is continuous if for all open sets $V\subseteq M$ the preimage $\gamma^{-1}(V)$ is open in the domain: on the domain we need
a topology; stillschweigend, when nothing else is said, $\R^d$ (and $\R$) is equipped with the
\emph{standard topology}, to connect with our idea that physical space is locally like $\R^d$ also as far as the topology
is concerned \ts{3}{1}{8:12}.
The same for a function $f\colon M\to\R$ (a temperature distribution on the manifold): continuous if preimages of open
sets are open; nothing more to say \ts{3}{1}{9:59}. If we only dealt with topological
manifolds and continuity, and not velocity vectors, we would be done: we would need no charts and no atlas
\ts{3}{1}{10:40}.

\subsection{View B: the chart level}\label{ssec:dm-view-b}

Recall that a chart is a pair $(U,x)$: $U\subseteq M$ open and $x\colon U\to x(U)\subseteq\R^d$ a map that is
continuous, invertible, with continuous inverse \ts{3}{1}{11:01}.
Everything seen in the chart has a one-to-one counterpart in the ``real world'': one would not draw a map of France
with the Eiffel tower twice; and if something is in a neighbourhood on the map, it is in a neighbourhood in the world: that
is continuity \ts{3}{1}{13:27}.

We have just decided continuity of a curve on the manifold, \emph{independently of the choice of chart}; the curve may even run
outside a chosen chart domain, which does not matter \ts{3}{1}{13:47}. Always think of the
manifold as the real world, what really happens, and the chart as our picture of it.

\begin{aside}[Blind spots of the eyes]
Told in class, referring to a problem on the new sheet \ts{3}{1}{14:29}:
each eye has a blind spot; close one eye and there is a point one does not see; the brain smooths this out. Each eye
maps the world onto a chart (the retina) with one point missing; the brain \emph{overlaps} the two charts and so builds a picture of
the real world with a lot of overlap. ``Seeing with two eyes has been exercised for 550 million years''
(the claim made in class, not checked by the notes).
\end{aside}

Take a point of the blue curve inside the chart domain; the chart image of the curve is the \emph{course drawn by the
captain on his sea chart}; with another sea chart the drawn course is different, yet the course in the real world, on the
level of the manifold, is the same \ts{3}{1}{15:51}. As physicists we think ``there is the real
world and something happens there'', but our \emph{description} of physics always happens in charts: charts give us real
numbers and coordinates, on which we can write formulas. The price is that a chart must be chosen, arbitrarily; so
our description in the chart must be made such that it is in the end \emph{independent of the chosen chart}. In school
and so far one always did physics only in the chart while thinking of describing the world; it is important to separate the two
\ts{3}{1}{16:33}.

\subsection{Commutative diagrams}\label{ssec:dm-diagrams}

Can we decide continuity of the curve also \emph{in the chart}? The first step: turn this ``kindergarten picture'' into a
``university picture'', i.e.\ write the maps and arrows with fewer prejudices; the technique is the \textbf{commutative
diagram} \ts{3}{1}{17:15}.

We have a curve $\gamma\colon\R\to M$ and a chart $(U,x)$ (here $U$ is open, so $x^{-1}$ continuous implies that the preimages of
open sets are open) \ts{3}{1}{18:22}. The curve does not care which charts one chooses: it runs
into the chart and out of it, or into another one \ts{3}{1}{19:29}. We choose a chart and concentrate on the part of the
curve inside the chart domain: the restriction
\[
  \gamma|_{\gamma^{-1}(U)}\colon\ \gamma^{-1}(U)\subseteq\R\to U ,
\]
whose domain is the preimage of $U$, i.e.\ the parameter values that land in $U$ (this restriction is necessary because for
some parameter values the curve lies outside the chart) \ts{3}{1}{19:49}.
Then apply the chart map:

\begin{definition}[New concepts: chart representation of a curve]\label{def:dm-chartrep}
For a curve $\gamma\colon\R\to M$ and a chart $(U,x)$, the \textbf{chart representation} (Kartendarstellung)
of $\gamma$ is the composition
\[
  x\circ\gamma\colon\ \gamma^{-1}(U)\subseteq\R\ \longrightarrow\ x(U)\subseteq\R^d,\qquad
  \lambda\mapsto\bigl(x^1(\gamma(\lambda)),\dots,x^d(\gamma(\lambda))\bigr)
\]
(the circle $\circ$ means ``first this, then that'': first the parameter into the real world, then the chart map down to the chart)
\ts{3}{1}{21:12}.
\end{definition}
This is the yellow piece of curve: the \emph{course drawn in the chart} (the ship's course on the chart). It is what the
physics courses wrote as $x^i(\lambda)$: the position coordinate of the particle at time $\lambda$
\ts{3}{1}{22:55}. This map, from a subset of $\R$ to $\R^d$, is called the
\emph{representation of $\gamma$ with respect to the chart $(U,x)$}; with another chart one in general gets another
representation, just as the course looks different on another sea chart \ts{3}{1}{23:56}.

\subsection{Atlas-continuity}\label{ssec:dm-atlas-cont}

The second view (view B) is: instead of looking at the curve in the real world, look at the curve \emph{in the chart}, which is
equipped with the standard topology of $\R^d$ and where we can decide continuity: the curve makes no jumps in the chart.
Can we conclude from this that it made none in the real world? That is the question \ts{3}{1}{24:38}.
The idea: to decide continuity of $\gamma$ on the manifold by investigating continuity of the \emph{chart representations}
\ts{3}{1}{25:19}.

\begin{definition}[New concepts: atlas-continuity]\label{def:dm-atlascont}
Let $(M,\mathcal O)$ be a topological manifold and $\mathcal A$ an atlas of $M$. A curve $\gamma\colon\R\to M$ is called
\textbf{$\mathcal A$-continuous} (continuous with respect to the atlas) if for every chart $(U,x)\in\mathcal A$ the chart
representation $x\circ\gamma|_{\gamma^{-1}(U)}\colon\gamma^{-1}(U)\to\R^d$ is continuous
\ts{3}{1}{26:06}.
\end{definition}
Here the map goes between subsets of $\R$ and $\R^d$, both with the standard topology (the domain with the topology induced
from $\R$, see the last problem sheet), so continuity is decided with school methods \ts{3}{1}{28:53}.
The promise: atlas-continuity turns out to be \emph{equivalent} to continuity, but a very important lesson will be
learned on the way \ts{3}{1}{26:27}. ``Since it depends on the choice of an atlas, the atlas is the collection of
charts whose chart domains cover the whole world''; the definition requires, in each chart image, a continuous piece
\ts{3}{1}{29:54}.

\begin{aside}[Notation for the structure]
Schuller writes sometimes simply ``topological manifold $M$'' and means the full structure $(M,\mathcal O,\mathcal A)$:
``at some point one has piled up so many structures that one calls the whole thing by its first name instead of its full
name; you must always complete this in your head. I try for a while to write it out so that you see the layers''
\ts{3}{1}{27:09}.
\end{aside}

\subsection{Why this has a chance: the lemma}\label{ssec:dm-lemma}

Remark~1: why does such an \emph{alternative definition} of continuity have any chance of agreeing with the usual one?
\ts{3}{1}{31:57}. The answer is a small mathematical lemma (Hilfssatz) \ts{3}{1}{33:02}:

\begin{theorem}[Lemma: continuity of a composition]\label{thm:dm-lemma}
Let $(A,\mathcal O_A)$, $(B,\mathcal O_B)$, $(C,\mathcal O_C)$ be topological spaces and let $f\colon B\to C$ be continuous.
Then a map $g\colon A\to B$ is continuous if and only if $f\circ g\colon A\to C$ is continuous
\ts{3}{1}{33:02}.
\end{theorem}
\begin{proof}
As a diagram: spaces $A\xrightarrow{g}B\xrightarrow{f}C$, with $f$ continuous by assumption and nothing assumed about $g$ or
$f\circ g$; the arrow name is written on top of the arrow \ts{3}{1}{35:11}. An equivalence is proved
in both directions \ts{3}{1}{35:52}. Continuity means: preimages of open sets are open. Let $W\subseteq C$ be
an arbitrary open set. Then
\begin{align*}
  (f\circ g)^{-1}(W)&=\{a\in A\mid f(g(a))\in W\}\\
  &=\{a\in A\mid g(a)\in f^{-1}(W)\}=g^{-1}\bigl(f^{-1}(W)\bigr),
\end{align*}
by applying the definition of the preimage twice \ts{3}{1}{36:13}.
\begin{itemize}
  \item ($\Rightarrow$) If $g$ is continuous: since $f$ is continuous (assumption), $f^{-1}(W)$ is open in $B$; since $g$ is
        continuous, $g^{-1}(f^{-1}(W))$ is open in $A$. This holds for every open $W$, so $f\circ g$ is continuous
        \ts{3}{1}{39:06}.
  \item ($\Leftarrow$) If $f\circ g$ is continuous, $(f\circ g)^{-1}(W)$ is open for all open $W\subseteq C$; this is
        $g^{-1}(f^{-1}(W))$, and it is open; so $g$ is continuous (as concluded in class; see the course note on the
        hypothesis on $f$) \ts{3}{1}{40:09}.
\end{itemize}
\end{proof}
\begin{coursenote}[The lemma needs $f$ to be a homeomorphism onto its image]
As stated in class (only $f$ continuous) the direction ($\Leftarrow$) is not valid in general: from the openness of
$g^{-1}(f^{-1}(W))$ for all open $W$ one cannot conclude that $g^{-1}(V)$ is open for \emph{every} open $V\subseteq B$
(take $f$ constant). It holds if $f$ is continuous, bijective onto its image with continuous inverse, because then
$g=f^{-1}\circ(f\circ g)$ is a composition of continuous maps. This is exactly the situation of the application below
($f=x$ is a chart map), so the use made in class is correct; the extra hypothesis is not stated in class
\ts{3}{1}{40:30}.
\end{coursenote}
\emph{Application.} The lemma is exactly the situation of the diagram above: $g$ is the curve restricted to the chart
domain, $f=x$ the chart map, which is continuous by definition of a chart; $f\circ g=x\circ\gamma$ the chart
representation. By the lemma, it is \emph{equivalent} whether we speak of continuity of the piece of curve in the real
world or of its image in the chart: ``this is the answer to why such a definition of continuity has a chance, indeed it
is actually equivalent'' \ts{3}{1}{40:53}.

\subsection{Chart independence and transition maps}\label{ssec:dm-transition}

A small but very important pedantic point \ts{3}{1}{42:17}: if we choose another atlas (an atlas is a figment of the mind; other
charts overlapping in other ways), we must ensure that we do not get a different definition. This is
\textbf{chart independence}, the ``physics without telekinesis'': we do not want things to move differently because
\emph{we} think another atlas \ts{3}{1}{42:17}. So Remark~2: is the definition of
$\mathcal A$-continuity \textbf{well defined} (wohldefiniert), i.e.\ independent of an additionally chosen structure?
\ts{3}{1}{43:18}

Let $(U,x)$ and $(V,y)$ be two charts, first of the same atlas. In principle, a piece of the curve could be continuous in
one chart and discontinuous in the other \ts{3}{1}{44:06}. By comparison with the continuity in the real
world it cannot, but suppose we had no access to this abstract definition: the captain in his cabin without windows has only
the atlas with certain transition functions, and must decide whether what is continuous in one chart is continuous in the
other \ts{3}{1}{44:26}. Two possibilities:
\begin{itemize}
  \item the chart domains $U$ and $V$ do not overlap: then the pieces of curve are different, cannot be compared, and there
        is no conflict; the charts are automatically ``compatible'' \ts{3}{1}{45:33};
  \item the chart domains overlap in $U\cap V\ne\emptyset$: look at the part $\gamma|_{\gamma^{-1}(U\cap V)}$ of the curve in
        the overlap; both $x$ and $y$ are defined there, so we have two representations of the same piece
        \ts{3}{1}{46:14}.
\end{itemize}
Go from the one chart image to the other: go back with $x^{-1}$ (which exists for a chart) and then down with $y$:
\[
  y\circ x^{-1}\colon\ x(U\cap V)\subseteq\R^d\ \longrightarrow\ y(U\cap V)\subseteq\R^d
\]
\ts{3}{1}{48:18}. This map goes directly between
(subsets of) $\R^d$ and $\R^d$, from one chart image to the other; in German one would say ``Blättern'' (``turning pages'').

\begin{definition}[New concepts: chart transition map]\label{def:dm-transition}
For two charts $(U,x)$ and $(V,y)$ of $M$ with $U\cap V\ne\emptyset$, the \textbf{chart transition map}
(Kartenwechselabbildung) from $x$ to $y$ is
\[
  y\circ x^{-1}\colon\ x(U\cap V)\to y(U\cap V),
\]
a map between open subsets of $\R^d$; likewise the map in the other direction $x\circ y^{-1}$
\ts{3}{1}{55:24}.
\end{definition}
Apply the lemma to the representations $x\circ\gamma|_{U\cap V}$ and $y\circ\gamma|_{U\cap V}$: since the chart maps
are continuous (and invertible with continuous inverse), $x\circ\gamma|_{U\cap V}$ is continuous if and only if
$y\circ\gamma|_{U\cap V}$ is continuous: the chart-level continuity does not depend on the chart or the atlas
\ts{3}{1}{49:40}.

\begin{intuition}[Why the chart maps must be invertible with continuous inverse]
Now one understands the requirement that chart maps and their inverses are continuous: if chart maps and their inverses were
not continuous, the chart transition map would not be continuous, and in the composition (the way down from one chart,
up to the world and down to the other) the continuity would be destroyed. ``The continuity of the transition map is a
necessary and a sufficient condition for the continuity to be decidable on the chart level independently of the chart''
\ts{3}{1}{52:08}. If it were only about
continuity, we would not need this whole section; we looked at the chart level in order to arrive at the point that the
transition functions must \emph{preserve} the property we want to define in the chart \ts{3}{1}{53:51}.
\end{intuition}

\subsection{$C^k$-compatibility}\label{ssec:dm-compat}

\begin{definition}[New concepts: $C^k$-compatible charts, $C^k$ atlas]\label{def:dm-compat}
Two charts $(U,x)$ and $(V,y)$ are called \textbf{$C^k$-compatible} if either $U\cap V=\emptyset$, or $U\cap V\ne\emptyset$
and the transition maps $y\circ x^{-1}\colon x(U\cap V)\to y(U\cap V)$ and $x\circ y^{-1}$ are \textbf{$k$-times continuously
differentiable}, i.e.\ of class $C^k$ as maps between subsets of $\R^d$ (standard analysis from the first semesters)
\ts{3}{1}{56:59}.
An atlas $\mathcal A$ of a topological manifold is called a \textbf{$C^k$ atlas} if every pair of overlapping charts in it
is $C^k$-compatible \ts{3}{1}{58:48}.
\end{definition}
The plural ``transition maps'' is because there are two directions, both required to have the property; the
requirement is symmetric \ts{3}{1}{58:06}. (A two-charts case: the requirement is
needed in both directions.)

\begin{theorem}[Every atlas is a $C^0$ atlas]\label{thm:dm-c0}
Every atlas of a topological manifold, as defined last time, is automatically a $C^0$ atlas: all transition maps are
continuous, since they are compositions of the continuous maps $y$ and $x^{-1}$ (read off the lower triangle of the diagram)
\ts{3}{1}{1:00:16}.
\end{theorem}
For an atlas as defined so far we know nothing more than that the transition maps are at least continuous
\ts{3}{1}{1:01:40}. After the break the motivation follows to demand sharper conditions: a true $C^1$, $C^2$
or even $C^\infty$ atlas \ts{3}{1}{1:02:00}.

% ==================================================================
\section{Differentiable manifolds}\label{sec:dm-diffmfd}

\begin{aside}[Greetings]
Before the second half, Schuller recalls having suggested that students greet the professor and the professors merely
say ``good morning''; colleagues objected \ts{3}{1}{1:02:21}.
\end{aside}

By ``differentiable'' we mean that the manifold carries a $C^k$ atlas with $k>0$ \ts{3}{1}{1:02:55}.

\begin{definition}[New concepts: maximal $C^k$ atlas, $C^k$ manifold]\label{def:dm-ckmanifold}
A $C^k$ atlas $\mathcal A$ of a topological manifold $(M,\mathcal O)$ is called \textbf{maximal} (maximaler $C^k$-Atlas) if every
chart $(U,x)$ that is $C^k$-compatible with \emph{all} charts of $\mathcal A$ already belongs to $\mathcal A$
\ts{3}{1}{1:04:25}.
A topological manifold together with a (maximal) $C^k$ atlas, $(M,\mathcal O,\mathcal A)$, is called a
\textbf{$C^k$ manifold} (differentiable manifold) \ts{3}{1}{1:03:38}.
\end{definition}
For $k=0$ one does not need to name the atlas, since any atlas is $C^0$ and one can reconstruct it; but a $C^k$ atlas with
$k>0$ is extra structure that must be given \ts{3}{1}{1:04:03}.

\begin{intuition}[The maximal atlas as Christmas present with a photocopier]
Suppose you give me a world atlas, a nice $C^0$ atlas. I say: page 204 is lovely but if it were tilted a little, my own
café would be on a vertical with Paris; so I photocopy it with the slight tilt and bind it in. Someone else wants it shifted a little
in another direction: also fine, in principle pointless because you can already read everything, but \emph{perhaps someone
wants that extra chart}. The maximal atlas satisfies all such wishes: it contains every chart that satisfies the
compatibility condition with the existing ones \ts{3}{1}{1:05:53}.\par
We will never write such an atlas down explicitly; we will always give an ordinary atlas, and say ``it is a maximal atlas'':
you get the photocopier as well, which may distort in any $C^k$ way you want, and prints all the charts that are
compatible with the original ones \ts{3}{1}{1:07:16}. The maximal-atlas principle
allows afterwards everything that has $C^k$ transition functions: a typical mathematical story. A $C^k$ manifold is one
that is equipped with some concrete atlas that we may always think extended to the maximal atlas based on it
\ts{3}{1}{1:07:58}.
\end{intuition}

A topological manifold always carries a maximal $C^0$ atlas. If one wants a $C^1$ atlas, one has to \emph{tear out charts} from
the maximal $C^0$ atlas, so that for the remaining ones all transition maps are not only continuous but once continuously
differentiable; for a $C^2$ atlas one must tear out still more. Which charts are torn out is a choice: one set of torn-out
charts leaves a $C^1$ atlas, another set another one \ts{3}{1}{1:08:40}.
Whenever one makes such a choice one gives the object \emph{more structure}: one could code secret messages in which pages
one tears out; just as before, with a set one chose which subsets are open (structure), here one gives structure to the
topological manifold \ts{3}{1}{1:10:03}.

\begin{intuition}[More structure, more theorems]
The more structure we pile up, the more special the choices, and the more can be done: ``the most powerful mathematics is that
which makes no assumptions'' sounds good but then nothing can be proved; with many assumptions one can prove a lot, but it
is very special. This is always the trade-off \ts{3}{1}{1:11:07}.
\end{intuition}

\subsection{Differentiable functions}\label{ssec:dm-diff-fn}

What can one do with a $C^k$ manifold? The answer is: define \emph{differentiability} of curves, functions, maps
\ts{3}{1}{1:11:27}. Schuller does it for functions now (for curves it works the same)
\ts{3}{1}{1:12:11}.

Let $f\colon M\to\R$ be a function and $(M,\mathcal O,\mathcal A)$ a $C^k$ manifold. On the manifold level one cannot even
think of deciding differentiability: it is merely a set with a topology; we must go down to the charts, ``the dirty
chart work'' \ts{3}{1}{1:12:32}.
Take a chart $(U,x)\in\mathcal A$ and consider the restriction $f|_U$; the chart is invertible, so
we can go up with $x^{-1}$ and then apply $f$:

\begin{definition}[New concepts: chart representation of a function]\label{def:dm-fnrep}
The \textbf{chart representation} of $f\colon M\to\R$ with respect to $(U,x)$ is
\[
  f\circ x^{-1}\colon\ x(U)\subseteq\R^d\ \longrightarrow\ \R
\]
(it is a function of $d$ real variables; the representation of the function in the chart)
\ts{3}{1}{1:13:34}.
\end{definition}
\begin{intuition}[Temperature on the sea chart]
In the real world there is a temperature function on a region (for example a plate or physical space); visualise it by
\emph{level lines} (isotherms): connecting all points of value 5, all points of value 7, and so on. In the chart, the image of
the region carries the level lines of $f\circ x^{-1}$, which may look quite different (the blue lines in the chart versus
the red ones in the real world). The captain draws on his sea chart the ocean temperatures 15, 2, 0 degrees, whereas in the
world there is the really prevailing temperature distribution: it is exactly like curves, only for functions
\ts{3}{1}{1:15:57}.
\end{intuition}
One could define the function to be continuous if its chart representation is continuous; with a second chart one adds the
transition map, which is continuous again: nothing new, as for curves \ts{3}{1}{1:17:41}.
But we want more: the function in the real world is a function on the manifold, and ``I cannot really say when it is
differentiable; I must say when we \emph{want to call it} differentiable'': when its chart representation is. Then we see that
a $C^k$ manifold is required for the definition to be well defined \ts{3}{1}{1:18:02}.

\begin{definition}[New concepts: $C^m$ function on a $C^k$ manifold]\label{def:dm-cm-fn}
Let $(M,\mathcal O,\mathcal A)$ be a $C^k$ manifold. A function $f\colon M\to\R$ is called \textbf{$m$-times continuously
differentiable} ($C^m$) if for every chart $(U,x)\in\mathcal A$ the chart representation
$f\circ x^{-1}\colon x(U)\subseteq\R^d\to\R$ is $m$-times continuously differentiable as an ordinary function of $d$ real
variables (analysis, first semesters) \ts{3}{1}{1:18:43}.
\end{definition}
In contrast to continuity, here there is nothing to compare with in the real world (continuity could be defined without
reference to the chart); the definition can only be made in the chart, so \emph{well-definedness} must be checked with special
care \ts{3}{1}{1:20:53}.

\subsection{Well-definedness: $k\ge m$}\label{ssec:dm-welldef}

\begin{theorem}[Well-definedness of $C^m$ for functions]\label{thm:dm-welldef}
Let $(M,\mathcal O,\mathcal A)$ be a $C^k$ manifold and $f\colon M\to\R$. If $f\circ x^{-1}$ is $C^m$ for one chart
$(U,x)\in\mathcal A$, then $f\circ y^{-1}$ is $C^m$ for every other chart $(V,y)\in\mathcal A$ on $U\cap V$, provided
$k\ge m$ \ts{3}{1}{1:21:34}.
\end{theorem}
\begin{proof}
Let $(U,x)$ and $(V,y)$ be two charts of $\mathcal A$ with $U\cap V\ne\emptyset$ (it is the same picture again and again,
``I shall ask you at some point to carry through this same number for curves, functions, maps, fields: it must be in
your flesh and blood''). There are two representations of the same objective thing in the real world, one in each of two atlases
(``once in the Westermann, once in the Diercke atlas'') \ts{3}{1}{1:21:54}.
Suppose we know that $f\circ x^{-1}$ is $C^m$ on $x(U\cap V)$; we want to conclude the same for $f\circ y^{-1}$
\ts{3}{1}{1:23:44}. Algebraically: insert the identity $x^{-1}\circ x$ (``it does nothing''):
\[
  f\circ y^{-1}\;=\;f\circ\bigl(x^{-1}\circ x\bigr)\circ y^{-1}\;=\;\bigl(f\circ x^{-1}\bigr)\circ\bigl(x\circ y^{-1}\bigr),
\]
using that the composition is associative \ts{3}{1}{1:24:26}.
(Schuller first wrote the right-hand side with the order of $x$ and $y$ interchanged and corrected himself with the help of a
student; the right order is the one above, and in the diagram it is: go from the chart $y$ up and down to the chart $x$, then
apply $f\circ x^{-1}$ \ts{3}{1}{1:28:00}.)
The first factor $f\circ x^{-1}$ is $C^m$ by assumption; the second, $x\circ y^{-1}$, is the transition map
between two charts of a $C^k$ atlas, hence $C^k$ \ts{3}{1}{1:25:27}.
The composition of a $C^k$ map and a $C^m$ map is $C^{\min(k,m)}$ \ts{3}{1}{1:26:30}. So
$f\circ y^{-1}\in C^{\min(k,m)}$, and this is $C^m$ if and only if $k\ge m$ \ts{3}{1}{1:26:51}.
\end{proof}
So to ensure that differentiability of a function is the same in every chart, one needs a $C^k$ atlas with $k$ at least
the desired differentiability class $m$ of the function: ``if I want to define $m$-times differentiable curves or
functions, I need charts whose transition functions preserve this property'' \ts{3}{1}{1:27:18}.
If we had only a $C^0$ atlas, a function could look differentiable in one chart and not in another; to avoid this one must
tear out enough pages from the atlas and give only $C^1$, $C^2$, $C^3$ atlases \ts{3}{1}{1:29:07}.

\begin{aside}[Why the formalism is the minimal version]
``You may sometimes be overwhelmed by the formalism, but the experience is that if I tell you the same things only in words
you leave the lecture saying it was all clear, and when I give you an exercise nothing is clear any more. This formalism is
the minimal, but complete, version of what you need; you can go on with general relativity in the master with it: it is the
most modern state of what mathematics has to offer''
\ts{3}{1}{1:29:50}. Take care, one does not get every detail at once: understand the diagrams, know
\emph{from where to where each map goes}: ``I must be able to wake you at night and you say the chart representation goes
from there to there''. If you do not know from where to where a map goes, you have not understood something in physics;
and if you do not draw the diagrams you never will \ts{3}{1}{1:31:00}.
\end{aside}

\subsection{Terminology}\label{ssec:dm-terminology}

A $C^\infty$ manifold is also called a \textbf{smooth manifold} (glatte Mannigfaltigkeit) \ts{3}{1}{1:31:20}.
The term ``differentiable manifold'' is actually a bad terminology, used only because it stands everywhere: it is not
the manifold that is differentiable but the \emph{objects living on it} (curves, functions, \dots). Correctly it should be
``a manifold with an atlas that allows definitions of differentiability'', and in the end we are interested not only in
differentiability but in the actual differentiating, which comes next lecture; similarly for a ``differentiable atlas''
\ts{3}{1}{1:32:03}.

% ==================================================================
\section{A worked example}\label{sec:dm-example}

A simple example that clarifies possible misunderstandings; the concepts must be kept apart: \emph{chart domain, image of
the chart domain, chart map, function, representation of the function in a chart, curve, representation of the curve in
a chart}: all differently defined, and one must learn to distinguish them \ts{3}{1}{1:34:23}.
Keep in mind the final goal: understand space (spacetime) as a differentiable manifold, which then needs even more
structure; these examples are simple models of a space or spacetime to which structure is added layer by layer
\ts{3}{1}{1:35:04}.

\begin{example}[$\R$ with two charts]\label{ex:dm-line}
Let $M=\R$ (a one-dimensional manifold) with the standard topology: a topological manifold of dimension $1$, locally (and
globally) $\R^1$ \ts{3}{1}{1:35:45}. Take two charts with the same chart
domain $U=V=\R$ (so they cover everything):
\begin{align*}
  x\colon\R\to\R,&\qquad x(m)=m \qquad\text{(the identity)},\\
  y\colon\R\to\R,&\qquad y(m)=m^3 .
\end{align*}
Both are charts: they are continuous, invertible ($x^{-1}(m)=m$, $y^{-1}(m)=\sqrt[3]{m}$) and the inverses are continuous
(the cube root is continuous: ``it goes up like this, no jumps'') \ts{3}{1}{1:37:07}.
\begin{itemize}
  \item \emph{Atlas~1} $=\{(\R,x)\}$: the only transition map is $x\circ x^{-1}=\mathrm{id}$, which is $C^\infty$: a $C^\infty$
        atlas, so $\R$ with it is a smooth manifold \ts{3}{1}{1:39:10}.
  \item \emph{Atlas~2} $=\{(\R,y)\}$: likewise, only the identity as transition map, again a $C^\infty$ atlas
        \ts{3}{1}{1:40:14}.
  \item \emph{Union} $=\{(\R,x),(\R,y)\}$ (``Susi got atlas~1, Kurti atlas~2; glue the two atlases together, what can
        happen? it cannot hurt''): the transition maps are
        \[
          y\circ x^{-1}\colon m\mapsto m^3\quad(C^\infty),\qquad x\circ y^{-1}\colon m\mapsto\sqrt[3]{m}.
        \]
        The derivative of $\sqrt[3]{m}$ is $\tfrac13 m^{-2/3}$, which diverges at $m=0$: the cube root is continuous
        but \emph{not} differentiable at $0$. So the union is only a $C^0$ atlas, not even $C^1$
        \ts{3}{1}{1:40:55}.
\end{itemize}
\end{example}
The lesson: ``uniting atlases is dangerous'' \ts{3}{1}{1:43:26}. And one must not ask ``is the chart map $x$
differentiable?'': that is nonsense; it is only meaningful to ask whether the \emph{transition maps} are differentiable, after
one has assumed that it is clear what ``differentiable'' means (which is exactly what is to be defined)
\ts{3}{1}{1:43:47}. If one is given the union atlas ($C^0$), one must tear out charts to
get to a $C^\infty$ atlas: tearing out $y$ leaves the $C^\infty$ atlas~1, tearing out $x$ leaves the $C^\infty$
atlas~2; tearing out charts from a $C^0$ atlas yields a higher differentiability class
\ts{3}{1}{1:44:28}. Happy Easter, and see you on Thursday \ts{3}{1}{1:44:49}.

\begin{exercise}[Two smooth structures on $\R$]\label{ex:dm-two-structures}
With $x(m)=m$ and $y(m)=m^3$ as above, show that the two maximal $C^\infty$ atlases generated by $\{(\R,x)\}$ and
$\{(\R,y)\}$ are different (the chart $y$ is not $C^\infty$-compatible with $x$). Are the two manifolds nevertheless
diffeomorphic? (Added; not said in class. This is the standard example that \emph{different} smooth structures on the same
topological manifold may give diffeomorphic manifolds.)
\end{exercise}

% ==================================================================
\section*{Key concepts}
\begin{keyconcepts}
\textbf{Retrieval questions.}
\begin{enumerate}
  \item Why do velocity vectors require more structure than a topological manifold?
  \item Define the chart representation of a curve and of a function. Write the maps with domain and codomain.
  \item State the lemma on the continuity of a composition and use it to show that atlas-continuity is equivalent to continuity.
  \item What is a chart transition map? Why must chart maps and their inverses be continuous?
  \item When are two charts $C^k$-compatible? What is a $C^k$ atlas and a maximal atlas? What is a $C^k$ manifold?
  \item Why is $C^m$ of a function well defined only for $k\ge m$? Give the proof.
  \item Why is the union of the two atlases $\{x=\mathrm{id}\}$ and $\{y(m)=m^3\}$ on $\R$ only a $C^0$ atlas?
  \item Why is ``differentiable manifold'' bad terminology?
\end{enumerate}
\textbf{Concepts.}
\begin{itemize}
  \item Manifold level versus chart level; physics must be chart independent (\cref{sec:dm-two-views}).
  \item Chart representation $x\circ\gamma$, $f\circ x^{-1}$ (\cref{def:dm-chartrep,def:dm-fnrep}); atlas-continuity
        (\cref{def:dm-atlascont}); lemma (\cref{thm:dm-lemma}).
  \item Transition maps $y\circ x^{-1}$ (\cref{def:dm-transition}); $C^k$-compatibility and $C^k$ atlas
        (\cref{def:dm-compat}); every atlas is $C^0$ (\cref{thm:dm-c0}).
  \item Maximal atlas, $C^k$ manifold, smooth manifold (\cref{def:dm-ckmanifold}); $C^m$ functions and well-definedness
        for $k\ge m$ (\cref{def:dm-cm-fn,thm:dm-welldef}); the example $\R$ with $x$ and $y=m^3$ (\cref{ex:dm-line}).
\end{itemize}
\end{keyconcepts}

\flushside
\end{paracol}
